\section{Pipeline thị giác máy tính và lập kế hoạch}

\subsection{Vị trí của AI trong hệ thống}

AI trong hệ thống không trực tiếp điều khiển đèn. Vai trò của AI là chuyển ảnh giao thông thành dữ liệu có cấu trúc: phương tiện nào xuất hiện, ở hướng nào, thuộc lớp nào và tạo ra mật độ tương ứng ra sao. Dữ liệu này đi qua bộ lập kế hoạch để tạo thời gian xanh; ESP32 mới là nơi thực thi chu kỳ đèn. Vì vậy phần thị giác máy tính được trình bày sau phần nhúng, nhưng vẫn là nguồn dữ liệu quan trọng của hệ thống thích nghi.

\begin{figure}[H]
  \centering
  \safeincludegraphics[width=0.88\linewidth]{figures/atlc_five_stage_pipeline.png}
  \caption{Pipeline nhiều tầng từ phát hiện phương tiện, ROI/counting, density, scheduler đến thực thi.}
  \label{fig:five-stage}
\end{figure}

\subsection{Backend detector và chuẩn hóa output}

Module \texttt{detector\_factory} chọn backend từ cấu hình \texttt{DETECTOR\_BACKEND}. Khi dùng backend \texttt{yolo}, lớp \texttt{LocalYoloDetector} gọi Ultralytics YOLO cục bộ và che giấu chi tiết output của thư viện \cite{ultralytics_predict}. Các module phía sau chỉ nhận một danh sách detection có dạng thống nhất:

\begin{lstlisting}
{
  "bbox": [x1, y1, x2, y2],
  "conf": 0.87,
  "class_name": "car"
}
\end{lstlisting}

Thiết kế này giúp thay model mà không phải sửa ROI, density, scheduler hoặc firmware. Trong tương lai, model \texttt{.pt} có thể được thay bằng ONNX hoặc NCNN nếu cần tối ưu trên Raspberry Pi, miễn là output cuối cùng vẫn giữ schema detection nêu trên.

\placeholderfigure{figures/bounding_box.png}{Ví dụ bounding box, nhãn lớp và confidence score của detector YOLO.}{fig:bounding-box}

\subsection{Tự huấn luyện detector cho miền giao thông}

Một điểm quan trọng của dự án là không dừng ở việc dùng nguyên model YOLO huấn luyện sẵn. Ghi chú benchmark của repo cho thấy model tổng quát phát hiện ô tô, xe buýt và xe tải tương đối ổn, nhưng bỏ sót nhiều xe máy trong góc camera giao thông: xe máy nhỏ, bị che khuất, xuất hiện dày và khác miền dữ liệu huấn luyện phổ thông. Giảm ngưỡng tin cậy không giải quyết gốc vấn đề; tăng kích thước ảnh giúp một phần nhưng làm FPS giảm. Vì vậy repo có nhánh nghiên cứu \texttt{yolo\_research} để huấn luyện detector tùy biến cho bốn lớp phương tiện của bài toán.

\begin{table}[H]
\centering
\caption{Cấu hình huấn luyện YOLO tùy biến trong repo.}
\label{tab:yolo-training-config}
\small
\begin{tabularx}{\linewidth}{>{\bfseries}p{4.0cm}Y}
\toprule
Hạng mục & Giá trị \\
\midrule
Dataset & \texttt{/content/atlc\_2000} \\
Số ảnh/nhãn & 1.984 ảnh và 1.984 file nhãn \\
Số object hợp lệ & 15.658 object \\
Lớp & \texttt{car}, \texttt{motorbike}, \texttt{truck}, \texttt{bus} \\
Model gốc & \texttt{yolo26n.pt} \\
Cấu hình train & 50 epoch, ảnh 640 px, batch 16, GPU \texttt{device=0} \\
Ngưỡng inference đề xuất & \texttt{CONFIDENCE\_THRESHOLD=0.25} \\
Đường dẫn runtime đề xuất & \texttt{pc\_app/models/local/atlc\_yolo26n\_custom.pt} \\
\bottomrule
\end{tabularx}
\normalsize
\end{table}

Kết quả cuối của log huấn luyện \texttt{atlc\_yolo26n\_custom} ghi nhận precision khoảng 0.820, recall khoảng 0.808 và mAP50 khoảng 0.871 trên tập validation. Các số này không nên tách khỏi bối cảnh dataset và cấu hình huấn luyện, nhưng chúng cho thấy hướng tiếp cận đúng: cải thiện độ phù hợp miền dữ liệu trước khi bàn tới tối ưu FPS. Với mục tiêu của báo cáo Vi xử lý, detector tùy biến được xem là nguồn cảm nhận tốt hơn cho bộ điều khiển, không phải trọng tâm duy nhất của đề tài.

\subsection{ROI và đếm theo hướng}

ROI được lưu trong JSON, gồm polygon cho hướng A và hướng B. Với mỗi detection, hệ thống lấy tâm bounding box rồi dùng \texttt{cv2.pointPolygonTest} để kiểm tra tâm đó nằm trong ROI nào. Nếu tâm nằm trong cả hai ROI do vùng chồng lấn, hệ thống chọn ROI có tâm gần hơn theo trục \(x\). Nếu không nằm trong ROI nào, detection được đưa vào nhóm ngoài vùng quan tâm.

\placeholderfigure{figures/roi_example.png}{Minh họa vùng quan tâm dùng để tách phương tiện theo hướng đường.}{fig:roi}

Cách làm này có ba ưu điểm kỹ thuật. Thứ nhất, khi đổi camera hoặc video, người dùng chỉ cần tạo lại file ROI. Thứ hai, ROI tách khỏi detector, nên thay model không ảnh hưởng đến logic chia hướng. Thứ ba, kết quả đếm có thể kiểm thử bằng dữ liệu giả: cùng một bbox và polygon phải luôn cho cùng một hướng.

\subsection{Mật độ thô, PCE và EMA}

Sau khi chia hướng, pipeline tính ba mức tín hiệu:

\begin{table}[H]
\centering
\caption{Ba mức tín hiệu mật độ trong pipeline.}
\label{tab:density-levels}
\small
\begin{tabularx}{\linewidth}{>{\bfseries}p{3.4cm}YY}
\toprule
Tín hiệu & Công thức/nguồn & Ý nghĩa \\
\midrule
Mật độ thô & Số detection trong ROI & Dễ hiểu, dùng làm baseline. \\
Mật độ PCE & \(\sum n_{i,c}w_c\) & Tính đến ảnh hưởng khác nhau của từng loại xe. \\
Mật độ EMA & \(\alpha D_t + (1-\alpha)\hat{D}_{t-1}\) & Làm mượt nhiễu detection từng frame. \\
\bottomrule
\end{tabularx}
\normalsize
\end{table}

Trong công thức PCE, \(n_{i,c}\) là số xe lớp \(c\) ở hướng \(i\), còn \(w_c\) là trọng số tương đương xe con. Cấu hình hiện tại dùng motorbike 0.5, car 1.0, rickshaw 1.2 và bus/truck 2.5. Cách dùng PCE phù hợp với hướng nghiên cứu điều khiển đèn thích nghi có xét loại phương tiện, vì một xe buýt hoặc xe tải chiếm dụng không gian giao thông lớn hơn xe máy \cite{raza2025edge,nguyen_yolov8_density}.

EMA dùng \(\alpha=0.30\). Nếu \(\alpha\) quá cao, kế hoạch xanh dao động mạnh theo lỗi detector ngắn hạn. Nếu \(\alpha\) quá thấp, hệ thống phản ứng chậm trước thay đổi mật độ. Giá trị 0.30 là lựa chọn cân bằng cho prototype và cần được đánh giá lại khi chạy camera thật.

\subsection{Bộ lập kế hoạch dựa trên mật độ}

Scheduler hiện tại là rule-based để dễ giải thích và dễ kiểm thử. Hệ thống tính tỷ lệ mật độ:

\begin{equation}
r_A = \frac{\hat{D}_A}{\hat{D}_A + \hat{D}_B + \epsilon}, \quad
r_B = \frac{\hat{D}_B}{\hat{D}_A + \hat{D}_B + \epsilon}
\end{equation}

Tổng ngân sách xanh là \(2G_{base}\), sau đó chia theo tỷ lệ \(r_A\), \(r_B\) và clamp vào miền an toàn:

\begin{equation}
G_A = clamp(2G_{base}r_A, G_{min}, G_{max})
\end{equation}

\begin{equation}
G_B = clamp(2G_{base}r_B, G_{min}, G_{max})
\end{equation}

Với cấu hình mặc định, \(G_{base}=20\), \(G_{min}=10\), \(G_{max}=45\), vàng 3 giây và all-red 1 giây. Thiết kế này cố ý giữ vàng và all-red cố định trong nguyên mẫu để không mở rộng không cần thiết phần an toàn chuyển pha. Phần thích nghi chỉ nằm ở thời gian xanh.

\subsection{Giảm tải suy luận}

Pipeline có tham số \texttt{DETECT\_EVERY\_N\_FRAMES}. Khi \(N=1\), detector chạy trên mọi frame. Khi \(N=2\) hoặc \(N=3\), chương trình vẫn đọc video liên tục nhưng chỉ chạy YOLO theo chu kỳ, các frame còn lại dùng detection gần nhất. Đây là cách giảm tải đơn giản trên thiết bị biên. Đổi lại, nếu giao thông thay đổi rất nhanh, detection được tái sử dụng có thể trễ hơn thực tế. Do đó tham số này phải được benchmark theo phần cứng, nguồn video và model cụ thể \cite{opencv_videocapture}.
